% 电路与电子学III 课程笔记 —— 第2章（教材：孙立山《电路理论基础》第四版 第2章）
\chapter{线性直流电路}

\noindent
由直流电压源和（或）直流电流源激励，且各部分电压、电流都是恒定量的电路称为\textbf{直流电路}；由线性元件和独立源组成的直流电路称为\textbf{线性直流电路}，它是最简单、最基本的一类电路。本章内容分三部分：（1）\textbf{等效化简}——等效的概念、电阻串并联与星形三角形变换、含源支路的等效变换；（2）\textbf{一般分析方法}——支路电流法、回路电流法、节点电压法；（3）\textbf{运算放大器}及其电路的分析。

\section{电阻网络的等效}
若干电阻相互连接构成的电路称为\textbf{电阻网络}。所谓\textbf{等效}，是指被化简的网络与等效网络具有相同的对应端口特性，即端口 $u-i$ 关系方程相同。用等效网络代替原网络后，\textbf{不改变其余部分}（外电路）的电压和电流。

\corepoint{“等效”只对外电路成立，对内不等效：等效前后被替换部分内部的电压、电流、功率一般都会改变（例如电压源与电阻并联等效成电压源后，电阻上的电流与功率都不一样了）。这是判断“能不能这样等效”的第一原则。}

\subsection{电阻的串联}
当各电阻首尾依次连接、连接点处没有分支，由 KCL 可知各电阻流过相同电流，则称这些电阻为\textbf{串联}。由 KVL 及欧姆定律
\begin{equation}
u=R_1i+R_2i=(R_1+R_2)i=R_{eq}i
\end{equation}
故两个电阻串联的等效电阻为 $R_{eq}=R_1+R_2$；$n$ 个电阻串联时
\keyformula{R_{eq}=\sum_{k=1}^{n}R_k}
串联常用于\textbf{分压}，各电阻上的电压为
\keyformula{u_k=\frac{R_k}{R_{eq}}\,u}
两个电阻串联时分压公式为 $u_1=\dfrac{R_1}{R_1+R_2}u$，$u_2=\dfrac{R_2}{R_1+R_2}u$。串联时流过相同电流，各电阻消耗的功率 $p_k=R_ki^2$，故
\begin{equation}
\frac{u_1}{u_2}=\frac{p_1}{p_2}=\frac{R_1}{R_2}
\end{equation}
即\textbf{串联时各电阻上的电压和功率的分配均与电阻成正比}。

注意：某些情况下电阻并非依次连接，但只要通过相同的电流，也可按串联化简。但这种变换后若把参考点选在被变换部分内部，未变换部分的节点电压会改变，严格说来不能称为等效变换。

\subsection{电阻的并联}
若各电阻都接到同一对节点之间、承受相同电压，则称这些电阻为\textbf{并联}。由 KCL 及欧姆定律
\begin{equation}
i=\frac{u}{R_1}+\frac{u}{R_2}=(G_1+G_2)u=G_{eq}u
\end{equation}
故两个电阻并联的等效电导为 $G_{eq}=G_1+G_2$，等效电阻为 $R_{eq}=\dfrac{R_1R_2}{R_1+R_2}$；$n$ 个电阻并联时
\keyformula{G_{eq}=\sum_{k=1}^{n}G_k,
\qquad
\frac{1}{R_{eq}}=\sum_{k=1}^{n}\frac{1}{R_k}}
并联常用于\textbf{分流}，各电阻分得的电流为
\keyformula{i_k=\frac{G_k}{G_{eq}}\,i}
两个电阻并联时分流公式为 $i_1=\dfrac{R_2}{R_1+R_2}i$，$i_2=\dfrac{R_1}{R_1+R_2}i$。并联时承受相同电压，$p_k=G_ku^2$，故
\begin{equation}
\frac{i_1}{i_2}=\frac{p_1}{p_2}=\frac{G_1}{G_2}
\end{equation}
即\textbf{并联时各电阻上的电流和功率的分配均与电导成正比}。

\corepoint{串联看电阻（分压、功率与 $R$ 成正比），并联看电导（分流、功率与 $G$ 成正比）。两个电阻并联时“积除和”$R_1R_2/(R_1+R_2)$ 要背熟；$n$ 个相等电阻 $R$ 并联的等效电阻为 $R/n$。}

\subsection{电阻的星形联结和三角形联结}
电阻的星形联结简称 $Y$ 形联结或 $T$ 形联结；三角形联结简称 $\Delta$ 形联结或 $\pi$ 形联结。在一定条件下二者可以相互等效，用以实现网络的等效化简——典型场合是\textbf{电桥电路}，它不能用串并联等效直接计算等效电阻。

$Y$ 形与 $\Delta$ 形都是\textbf{三端网络}：有三对端子间电压 $u_{12}$、$u_{23}$、$u_{31}$ 和三个端子电流 $i_1$、$i_2$、$i_3$。由 KCL、KVL 恒有 $i_1+i_2+i_3=0$、$u_{12}+u_{23}+u_{31}=0$，故三端网络的对外作用可用两对端子间电压与两个端子电流的关系表示。令两种联结的对应端口特性方程系数相等，即得等效条件。

$\Delta$ 形联结到 $Y$ 形联结的等效变换公式为
\keyformula{R_1=\frac{R_{12}R_{31}}{R_{12}+R_{23}+R_{31}},\quad
R_2=\frac{R_{23}R_{12}}{R_{12}+R_{23}+R_{31}},\quad
R_3=\frac{R_{31}R_{23}}{R_{12}+R_{23}+R_{31}}}
可归纳为
\begin{equation}
R_{1,2,3}=\frac{\Delta\text{形联结中与该节点相连两电阻之积}}{\Delta\text{形联结中各电阻之和}}
\end{equation}
$Y$ 形联结到 $\Delta$ 形联结的等效变换公式为
\keyformula{R_{12}=\frac{R_1R_2+R_2R_3+R_3R_1}{R_3},\quad
R_{23}=\frac{R_1R_2+R_2R_3+R_3R_1}{R_1},\quad
R_{31}=\frac{R_1R_2+R_2R_3+R_3R_1}{R_2}}
可归纳为
\begin{equation}
R_{12,23,31}=\frac{Y\text{形联结各电阻两两乘积之和}}{Y\text{形联结中不相邻的电阻}}
\end{equation}

对称情况最简单常用：设三个相等电阻接成 $Y$ 形或 $\Delta$ 形，则
\keyformula{R_{\Delta}=3R_{Y},\qquad R_{Y}=\frac{1}{3}R_{\Delta}}
以上分析表明：对称 $Y$ 形联结等效为 $\Delta$ 形联结时也是对称的，反之亦然。

\corepoint{$Y$-$\Delta$ 变换的三个电阻必须“有公共节点”（$Y$）或“构成回路”（$\Delta$）才能用；变换时先标清三个端子的对应关系，再代公式，端子标错是最常见的错误。对称情形 $R_\Delta=3R_Y$ 可以快速验算。}

\section{含源支路的等效变换}
两种典型的含独立源一端口电路，在端口电压、电流取\textbf{非关联}参考方向的条件下，其电压电流关系分别为
\keyformula{U=U_s-R_iI,\qquad I=I_s-G_iU}
当电路参数满足一定条件时，\textbf{电压源和电阻串联支路与电流源和电导并联支路可以相互等效}。将第二式改写为
\begin{equation}
U=\frac{1}{G_i}I_s-\frac{1}{G_i}I
\end{equation}
与第一式比较，得两电路的等效变换条件
\keyformula{U_s=\frac{I_s}{G_i},\qquad R_i=\frac{1}{G_i}}
即 $U_s=I_sR_i$，$I_s=U_s/R_i$。

物理意义：实际稳压电源工作时，随端口电流变化，其端口电压也略有变化，故用“电压源＋电阻串联”作模型——$U_s$ 是空载（输出电流为零）时的输出电压，$R_iI$ 是由负载电流引起的输出电压降，$R_i$ 称为电源的等效内电阻。实际稳流电源则用“电流源＋电导并联”作模型——$I_s$ 是空载（输出电压为零即短路）时的输出电流，$G_iU$ 是由负载电压引起的输出电流下降，$G_i$ 称为内导。因此实际电源都有这两种等效的电路模型。

几种特殊连接：
\begin{itemize}
  \item 电压源与电阻\textbf{并联}的一端口，其端口电压恒为 $U_s$、电流不确定，对外作用等效为电压源 $U_s$（电阻可去掉）；
  \item 电流源与电阻\textbf{串联}的一端口，其端口电流恒为 $I_s$，对外作用等效为电流源 $I_s$（电阻可去掉）；
  \item 两个电压源串联、两个电流源并联，可直接按代数和合并；两个相同电压源并联、两个相同电流源串联，等效为一个同值电源。
\end{itemize}

受控电压源与电阻的串联也可等效成受控电流源与同一电阻的并联，或反之，变换方法与含独立源的情况相似；\textbf{但在变换过程中应保证控制量的位置不变}。

\corepoint{理想电压源（$R_i=0$）与理想电流源（$G_i=0$）之间\textbf{不能}相互等效，因为零不能取倒数。凡涉及理想电源的等效变换，都要先问“内阻是不是零”。}

\section{支路电流法}
能用等效化简分析的电路通常是结构较简单、只求个别量的电路。当电路复杂以致难以化简，或待求变量较多时，就要依据电路变量的基本约束列写方程组求解。

设电路有 $b$ 条支路、$n$ 个节点，各独立源、受控源和电阻的参数均已知，要求确定各支路电流或电压。\textbf{支路电流法}就是以 $b$ 条未知的支路电流为待求量：

\begin{enumerate}
  \item 设定各支路电流的参考方向，并标在电路图上；
  \item 对 $n-1$ 个独立节点列写 KCL 方程；
  \item 对 $b-(n-1)$ 个独立回路列写 KVL 方程，其中电阻电压用支路电流表示；
  \item 联立 $b$ 个方程求解，再按需由支路电流求其他量。
\end{enumerate}

支路电流法是分析电路的一种基本方法，正确运用的关键在于所列 $b$ 个方程\textbf{彼此独立}：$n-1$ 个节点的 KCL 方程是一组独立方程；为列独立的 KVL 方程，可每选一个回路时\textbf{至少包括一条新支路}，从而在该回路的 KVL 方程中包含一个新的支路电压，使方程独立于已列出的方程。

含电流源电路的处理：把电流源支路的电流视为已知，列 KCL 时可将其直接移到等号右端；对含电流源的回路列 KVL 时，须把未知的电流源端电压列入方程（多一个未知量、多一个补充方程），或\textbf{适当选取回路使所取回路不包含电流源}，从而减少方程数目。

含受控源电路的处理：在支路电流方程中，受控源先按独立源处理，再用支路电流表示其控制电压（流），代入后消去控制量。

\corepoint{支路电流法的方程数等于支路数 $b$，未知量最多、计算量最大，但思路最直接、不易出错，适合作为检验其他方法的基准。含电流源时“选回路绕开它”能立刻少一个方程。}

\section{回路电流法}
支路电流法的待求量数目等于未知支路电流数，所需方程较多。若以假想的回路电流为待求量，则可大幅减少方程数。

\subsection{回路电流}
由于电流具有连续性，可以假设在每个独立回路中分别存在一个闭合流动的电流，这就是\textbf{回路电流}；以网孔为回路时也称为\textbf{网孔电流}。假设回路电流后，支路电流便可用回路电流的代数和表示：若支路电流与回路电流方向相同，回路电流前取正号，否则取负号。

以回路电流为待求量便\textbf{等价地应用了 KCL}——把支路电流用回路电流表示后代入任一节点的 KCL 方程，得到的是恒等于零的恒等式。因此回路电流法\textbf{无需再列 KCL 方程}。

\subsection{回路电流方程的列写}
选择 $b-(n-1)$ 个独立回路，以各回路电流为待求量列写 KVL 方程，这种方法称为\textbf{回路电流法}（回路分析法）。其标准形式为
\keyformula{\begin{aligned}
R_{11}I_{m1}+R_{12}I_{m2}+\cdots+R_{1m}I_{mm}&=U_{s11}\\
R_{21}I_{m1}+R_{22}I_{m2}+\cdots+R_{2m}I_{mm}&=U_{s22}\\
&\vdots\\
R_{m1}I_{m1}+R_{m2}I_{m2}+\cdots+R_{mm}I_{mm}&=U_{smm}
\end{aligned}}
各项意义与列写规则：
\begin{enumerate}
  \item \textbf{自阻} $R_{kk}$：组成第 $k$ 个回路的各支路电阻之和，恒取正号；
  \item \textbf{互阻} $R_{kj}$（$k\neq j$）：第 $k$、$j$ 两回路公共支路上的电阻。若两回路电流在该公共支路上方向\textbf{相同}，互阻取正；\textbf{相反}取负。在仅由独立源和二端电阻组成的电路中 $R_{kj}=R_{jk}$，即系数矩阵\textbf{对称}；
  \item 等号右边 $U_{skk}$：沿第 $k$ 个回路方向各电压源\textbf{电位升}的代数和（电位升取正、电位降取负）。
\end{enumerate}
写成矩阵形式 $\boldsymbol{RI}_m=\boldsymbol{U}_s$，其中 $\boldsymbol{I}_m$ 称为回路电流向量，$\boldsymbol{U}_s$ 称为回路源电压向量，系数矩阵称为回路电阻矩阵。

含电流源电路：方法一是以网孔为独立回路，把电流源两端未知电压 $U$ 列入方程，再补充电流源的特性方程；方法二是\textbf{适当选取独立回路使电流源只包含在一个回路中}，则该回路电流即等于源电流 $I_s$，从而减少一个待求量及方程数。后者显然更优。

含受控源电路：先用回路电流表示受控源的控制量，代入消去控制量后整理。此时受控源的控制系数要影响方程的系数，通常使系数矩阵\textbf{不再具有对称性}。

\corepoint{回路电流法的三个易错点：（1）自阻恒正、互阻看方向；（2）等号右边是“电位升”的代数和（沿回路方向从“$-$”到“$+$”取正）；（3）含电流源时要主动选回路，让电流源独占一个回路。}

\section{节点电压法}
在集中参数电路中任选一节点为\textbf{参考点}，其他各节点与参考点之间的电压称为该节点的\textbf{节点电压}（节点电位）；参考点的节点电压为零。有了节点电压，任意两点间电压便可表达成这两个节点电压之差。

用节点电压表示支路电压，相当于\textbf{等价地列写了 KVL 方程}——绕任一回路一周，各节点电压相互抵消，方程恒等于零。因此以节点电压为待求量时\textbf{无需再列 KVL 方程}。以 $n-1$ 个节点电压为待求量、对 $n-1$ 个节点列写 KCL 方程的方法，就是\textbf{节点电压法}（节点分析法）。

其标准形式为
\keyformula{\begin{aligned}
G_{11}U_{n1}+G_{12}U_{n2}+\cdots+G_{1,n-1}U_{n,n-1}&=I_{s11}\\
G_{21}U_{n1}+G_{22}U_{n2}+\cdots+G_{2,n-1}U_{n,n-1}&=I_{s22}\\
&\vdots
\end{aligned}}
列写规则：
\begin{enumerate}
  \item \textbf{自导} $G_{kk}$：与节点 $k$ 直接相连的各支路电导之和，恒取正号；
  \item \textbf{互导} $G_{kj}$（$k\neq j$）：直接连接在节点 $k$、$j$ 之间的诸支路电导之和并\textbf{带负号}。在仅由独立源和线性二端电阻组成的电路中 $G_{kj}=G_{jk}$，系数矩阵\textbf{对称}；
  \item 等号右边 \textbf{节点源电流} $I_{skk}$：包括两项——（1）与节点 $k$ 相连的各电流源电流的代数和（电流\textbf{流入}节点取正，流出取负）；（2）与节点 $k$ 相连的各“电压源与串联电导”的乘积 $G U_s$ 的代数和（电压源\textbf{正极性端指向}该节点时取正，否则取负）。
\end{enumerate}
写成矩阵形式 $\boldsymbol{GU}_n=\boldsymbol{I}_s$，$\boldsymbol{U}_n$ 为节点电压向量，$\boldsymbol{I}_s$ 为节点源电流向量，系数矩阵为节点电导矩阵。复杂电路中节点电导矩阵元素大部分为零，属于\textbf{稀疏矩阵}，计算时要用稀疏矩阵技术。

只有一个独立节点（多条支路并联）时，节点电压可直接写出
\begin{equation}
U_{n1}=\frac{\sum I_s+\sum GU_s}{\sum G}
\end{equation}
即“注入电流的代数和 ÷ 各支路电导之和”，这就是\textbf{弥尔曼定理}。

\subsection{改进节点电压法}
上述节点电压法是根据欧姆定律用节点电压表示二端电阻的电流。若电路中存在\textbf{非二端电阻元件}（如纯电压源支路），该支路电流不能用支路电导与节点电压之积表示。解决办法有两种：

\begin{enumerate}
  \item \textbf{改进节点法}：以 $n-1$ 个节点电压和纯电压源支路的电流为待求量，列 $n-1$ 个节点的 KCL 方程，再与纯电压源支路的方程（$U_{nk}-U_{nj}=U_s$）联立；
  \item 若电路中只有一个纯电压源支路，可\textbf{选电压源的一端为参考点}，则另一端的节点电压即为已知量，问题大为简化——此时只需对其余节点列 KCL，且未知的电压源电流只出现在被略去的那个方程中。
\end{enumerate}

\corepoint{三种方法对比：支路电流法 $b$ 个方程（最直观）；回路电流法 $b-n+1$ 个方程（自阻互阻规则稍繁）；节点电压法 $n-1$ 个方程（自导互导规则最简单，便于编程）。因为节点电压法的列写规则比回路电流法简单，工程与仿真软件（PSPICE 等）普遍采用节点电压法。}

\begin{center}
\begin{tabular}{@{}llll@{}}
\toprule
方法 & 待求量 & 方程数 & 方程类型 \\
\midrule
支路电流法 & $b$ 个支路电流 & $b$ & $n-1$ 个 KCL $+\,b-n+1$ 个 KVL \\
回路电流法 & $b-n+1$ 个回路电流 & $b-n+1$ & 仅 KVL \\
节点电压法 & $n-1$ 个节点电压 & $n-1$ & 仅 KCL \\
\bottomrule
\end{tabular}
\end{center}

\section{运算放大器}
运算放大器（简称运放）是用集成电路工艺制成的\textbf{多端元件}，其内部电路分析超出本课程范围。由于它具有对输入信号进行求和、放大、积分、微分等运算功能，故得此名。

\subsection{端口特性}
仅从外部端子分析，运放是一种高放大倍数的电压放大器。用 $u_+$、$u_-$ 分别表示同相输入端和反相输入端的节点电压，则输入端口电压即\textbf{差分输入电压}为
\begin{equation}
u_d=u_+-u_-
\end{equation}
在输出端开路情况下，输出电压 $u_o$ 与差分输入电压的关系曲线可用折线近似：水平部分称为运放的\textbf{饱和区}，中间部分（$-U_{cc}<u_d<+U_{cc}$）为\textbf{线性区}，在此区域内输出电压与输入电压成正比。二者之比称为\textbf{开环增益}（开环电压放大倍数）$A$，即 $u_o=Au_d$。

本书只讨论运放工作在\textbf{线性区}时的功能与分析方法。线性区内运放的电路模型由输入电阻 $R_i$、输出电阻 $R_o$ 和受控电压源 $Au_d$ 组成（含受控源，说明运放是有源元件）。典型参数：开环增益 $A\sim10^{5}$、输入电阻 $R_i\sim10^{6}\sim10^{13}\,\Omega$、输出电阻 $R_o\sim10\sim100\,\Omega$，工作电压 $\pm5\sim\pm24\,\mathrm{V}$。

由于开环增益非常大，一个微小的输入电压就足以使运放进入饱和区。因此为使运放工作在线性区，\textbf{必须引入负反馈}——将输出电压或输出电压的一部分回馈到反相输入端；若回馈到同相输入端则属正反馈。

\subsection{理想运算放大器}
仅分析电路工作原理时常用\textbf{理想运放}模型，其条件是：无穷大的开环增益、无穷大的输入电阻和零输出电阻。由此得到理想运放的端口特性
\keyformula{i_+=i_-=0,\qquad u_+=u_-}
其中 $i_+$、$i_-$ 为两个输入端的电流。这两个式子不含理想运放的输出电压 $u_o$ 和输出电流 $i_o$，它们将由外接电路决定。

第一式表明输入端口电流为零，类似于\textbf{断路}，称为\textbf{虚断}；第二式表明输入端口电压为零，又类似于\textbf{短路}，称为\textbf{虚短}。故可以说理想运放输入端口满足“虚断”和“虚短”条件。

\corepoint{“虚短、虚断”是分析含理想运放电路的唯一钥匙，但两者都\textbf{以负反馈为前提}：没有负反馈时运放工作在饱和区，$u_+\neq u_-$，虚短不成立。判断电路是否含负反馈，是做题的第一步。}

\section{含运算放大器电路的分析}
理想运放输入端口电压、电流均为零，而输出端口电压、电流取决于运放之外的电路，这些特性决定了分析方法的特殊性。

\subsection{反相放大器}
输入电压经 $R_1$ 接反相输入端，$R_f$ 从输出端接回反相输入端，同相输入端接地。由虚断 $i_+=i_-=0$ 得 $i_1=i_f$；由虚短 $u_-=u_+=0$（\textbf{虚地}）有 $\dfrac{u_1-0}{R_1}=\dfrac{0-u_o}{R_f}$，于是得
\keyformula{u_o=-\frac{R_f}{R_1}u_1}
可见输出电压与输入电压成正比、极性相反，故称反相放大器。其增益只与\textbf{外接电阻}有关，故可实现高精度、高稳定性的增益值。反相输入端电位近似为零但不接地，称为“虚地”。

\subsection{同相放大器}
输入电压加在同相输入端，反相输入端经 $R_1$ 接地、经 $R_f$ 接输出端。由虚短 $u_-=u_+=u_1$，再由虚断 $i_1=i_f$ 及分压关系得
\keyformula{u_o=\left(1+\frac{R_f}{R_1}\right)u_1}
它是增益大于 1 的电压控制电压源，输出电压与输入电压极性相同。若令 $R_1=\infty$、$R_f=0$，则 $u_o=u_1$，电路成为\textbf{电压跟随器}，常在多级电路中起隔离作用——它从输入电压源吸收的电流为零，即输入电压源不提供功率。

\subsection{加法器}
三个输入电压经 $R_1$、$R_2$、$R_3$ 接反相输入端，$R_f$ 接回输出端。由虚短得 $u_-=0$，由虚断得 $i_1+i_2+i_3=i_f$，于是
\keyformula{u_o=-\left(\frac{R_f}{R_1}u_1+\frac{R_f}{R_2}u_2+\frac{R_f}{R_3}u_3\right)}
即输出电压是输入电压的\textbf{线性组合}（加权和），各系数称为加权系数。若 $R_1=R_2=R_3=R_f$，则各加权系数相等为 $-1$，$u_o=-(u_1+u_2+u_3)$。

\subsection{差分放大器}
同时从两个输入端加入输入电压，便构成差分放大器。由虚短、虚断并利用分压公式与 KVL 可求得
\begin{equation}
u_o=\frac{R_2}{R_1}\left(1+\frac{R_1}{R_2}\right)\frac{R_4}{R_3+R_4}u_2-\frac{R_2}{R_1}u_1
\end{equation}
若外接电阻满足平衡条件 $\dfrac{R_2}{R_1}=\dfrac{R_4}{R_3}$，则化简为
\keyformula{u_o=\frac{R_2}{R_1}\left(u_2-u_1\right)}
即输出电压与两输入电压之差成正比。差分放大器的主要优点是：若两个输入电压因干扰受到\textbf{相同}的扰动（共模信号），则该扰动不会影响输出电压。

\subsection{用运放实现负电阻}
利用运放可以把正电阻变换为\textbf{等效负电阻}：按虚断、虚短列方程，可得端口等效电阻为负值（$R_{eq}<0$）。负电阻是有源元件，能对外发出功率，这与 §1.2 的结论一致。

\subsection{复杂含运放电路的分析}
上述只涉及简单电路。电路较复杂时需列写联立方程组，在各种方法中\textbf{节点电压法}是分析含运放电路的有效方法：把运放的两个输入端当作普通节点列 KCL（因虚断，输入端电流为零，两节点之间不存在互导），再补充理想运放的端口方程 $u_+=u_-$ 作为约束。运放的输出电流只出现在输出节点的 KCL 方程中，若不需要求它，可略去输出节点的 KCL 方程。

\corepoint{含运放电路的分析套路：（1）确认存在负反馈；（2）对输入端节点写 KCL（虚断使输入端不“吃掉”电流）；（3）补 $u_+=u_-$（虚短/虚地）；（4）解出输出与输入的关系。求输出电流时，输出节点的 KCL 方程不可略。}

\section{对偶原理}
（本节内容在教材中位于第 3 章 3.6 节，因与本章的回路法、节点法直接对偶，故提前介绍。）

如果两个电路的方程互为对偶，则称这两个电路\textbf{互为对偶}。把电路中的变量和参数按对偶关系替换，方程形式不变，这一规律称为\textbf{对偶原理}。常用对偶关系如下：

\begin{center}
\begin{tabular}{@{}llll@{}}
\toprule
电压 $u$ & 电流 $i$ & 电阻 $R$ & 电导 $G$ \\
电荷 $q$ & 磁链 $\Psi$ & 电容 $C$ & 电感 $L$ \\
电压源 $u_s$ & 电流源 $i_s$ & 开路 & 短路 \\
串联 & 并联 & 节点 & 网孔（回路） \\
KCL & KVL & 节点电压法 & 回路电流法 \\
$Y$ 形联结 & $\Delta$ 形联结 & 自导 & 自阻 \\
\bottomrule
\end{tabular}
\end{center}

应用要点：
\begin{enumerate}
  \item 对偶是\textbf{方程层面}的对应关系，两个互为对偶的电路其结构图互为“对偶图”——原电路的网孔对应对偶电路的节点，反之亦然；
  \item 已知一个电路的结论，按对偶关系逐项替换即得对偶电路的结论（例如由分压公式 $u_k=R_k/R_{eq}\cdot u$ 立即得到分流公式 $i_k=G_k/G_{eq}\cdot i$）；
  \item 对偶关系要求参考方向也作相应约定，否则会差一个负号。
\end{enumerate}

\corepoint{对偶原理是“少背一半公式”的工具：串并联、分压分流、$Y$-$\Delta$、回路法与节点法、$RL$ 与 $RC$ 暂态，全都可以两两对偶记忆。但要记住它只对方程与结构成立，不是物理上的等价。}

\endinput
